%===========================================================
% note-setup-box.tex
% Customized Setup with Obsidian Callouts & Categorized Packages
%===========================================================

% ==========================================
% 1. 页面排版与段落 (Page & Layout)颜色必须放在最前面
% ==========================================
\documentclass{book} % 必须改为 book 类，以支持 \frontmatter 和 \chapter

\usepackage[dvipsnames]{xcolor} % 确保颜色模型丰富
\usepackage{hyperref}
\hypersetup{
    colorlinks=true,
    linkcolor=,             % 保持默认或根据需要设定
    urlcolor=cyan
}

\usepackage[a4paper,top=2cm,bottom=2.5cm,left=2.5cm,right=2.5cm]{geometry}
\usepackage{indentfirst}    % 首行缩进
\usepackage{anyfontsize}    % 任意字体大小
\usepackage{appendix}       % 附录支持
\linespread{1.25}           % 全局行距 1.25
\usepackage{float}
\usepackage{titlesec}
\usepackage{fancyhdr}
\usepackage{csvsimple}
\usepackage{pagecolor}
\pagestyle{fancy}
\renewcommand{\headrulewidth}{0pt}
\fancyhead{}
\fancyfoot[C]{\thepage}

% ==========================================
% 2. 字体 (Fonts)
% ==========================================
\usepackage{type1cm}
\usepackage[UTF8]{ctex} 
\raggedbottom
%\usepackage{mathpazo}这个字体不错，有时间用用
%\usepackage{charter} % 这里的字体支持希腊字母花体
%\usepackage[expert]{mathdesign}
% ==========================================
% 3. 数学 (Math)
% ==========================================
\usepackage{amsmath,amssymb,amsfonts,amsthm}
\usepackage{mathtools}
\usepackage{mathrsfs}
\usepackage{bm}             % 粗体数学符号
\usepackage{extarrows}
\allowdisplaybreaks[4]
% 自定义数学算子快捷命令
\newcommand{\E}{\operatorname{E}}
\newcommand{\Var}{\operatorname{Var}}
\newcommand{\Prob}{\operatorname{P}}
\newcommand{\I}{\operatorname{I}}
\newcommand{\diff}{\mathrm{d}}

% ==========================================
% 4. 图形 (Graphics & TikZ)
% ==========================================
\usepackage{graphicx}
\usepackage{eso-pic}        % 添加背景图片支持
\graphicspath{
    {./figure/}{./figures/}{./image/}{./images/}{./graphic/}{./graphics/}{./picture/}{./pictures/}
}
\usepackage{subcaption}     % 子图支持
\usepackage{tikz}
\usetikzlibrary{decorations.pathmorphing} % TikZ 路径变形库

% ==========================================
% 5. 表格与列表 (Tables & Lists)
% ==========================================
\usepackage{array}
\usepackage{booktabs}       % 三线表 (toprule, midrule, bottomrule)
\usepackage{multirow,multicol}
\usepackage[shortlabels,inline]{enumitem} % 强大的列表环境
\usepackage{tabularx}
\usepackage{listings}



\definecolor{ObsTheorem}{RGB}{30, 150, 150}   % #1E9696 深青绿
\definecolor{ObsProp}{RGB}{80, 160, 100}      % #50A064 深绿
\definecolor{ObsLemma}{RGB}{110, 180, 60}       % 更纯粹的绿意
\definecolor{ObsCor}{RGB}{160, 174, 45}       % #A0AE2D 青柠绿
\definecolor{ObsClaim}{RGB}{210, 200, 35}     % 更明亮的黄色
\definecolor{ObsDef}{RGB}{110, 110, 245}      % #6E6EF5 偏蓝紫
\definecolor{ObsRemark}{RGB}{255, 140, 0}     % #FF8C00 橙色
\definecolor{ObsNote}{RGB}{255, 160, 180}     % #FFA0B4 浅粉
\definecolor{ObsSol}{RGB}{245, 170, 120}      % #F5AA78 暖橙
\definecolor{ObsEg}{RGB}{205, 170, 95}        % 更深沉的暗褐黄
\definecolor{ObsPro}{RGB}{150, 110, 210}      % 偏紫的色调
\definecolor{ObsEx}{HTML}{57A3F1}           % 原有的蓝%

% ==========================================
% 7. 代码与算法 (Code & Algorithms)
% ==========================================
\usepackage{minted}         % 现代代码高亮包 (需编译时加 -shell-escape)
% 算法包 (只保留您指定的一套格式)
\usepackage[linesnumbered,lined,boxed,commentsnumbered]{algorithm2e}
% ==========================================
% 8. 其它工具 (Other Utilities)
% ==========================================
\usepackage{cancel}         % 公式中的删除线
\usepackage{extarrows}      % 更多长箭头
\usepackage[normalem]{ulem} % 文本下划线与删除线
\usepackage{verbatim}       % 多行注释
\usepackage{float}          % 浮动体强行排版 (H)
\usepackage[most]{tcolorbox}% 彩色盒子核心宏包
\usepackage{fontawesome5}   % 【新增】FontAwesome 图标库，用于渲染 Obsidian 的 callout 图标！


% ==========================================
% 9 & 10. 定理环境样式定义与 Tcolorbox 绑定
% ==========================================
\renewcommand*{\proofname}{\normalfont\bfseries Proof}
\usepackage{thmtools}

% 核心宏：一键绑定 thmtools 样式与 tcolorbox 外观
\newcommand{\createobsbox}[4]{
    % 1. 配置专属 thmtools 样式
    \declaretheoremstyle[
        headfont=\bfseries\color{#3},      % 标题局部染色
        notefont=\bfseries\color{#3},      % 附加标题同色
        notebraces={}{},                   % 【修改】留空！去掉丑陋的自动外层圆括号
        bodyfont=\normalfont,              % 绝对不要在这里写 \color{black}！
        spaceabove=0pt, spacebelow=0pt,
        headpunct={},                      % 【修改】留空！去掉末尾多余的点号 (.)
        postheadspace=\newline             % 强制换行
    ]{#1style}

    % 2. 声明定理环境
    \declaretheorem[style=#1style, name=#2, #4]{#1}
    \declaretheorem[style=#1style, name=#2, numbered=no]{#1*}

    % 3. 绑定 tcolorbox 样式 (回归干净配置)
    \tcolorboxenvironment{#1}{
        standard, breakable, rounded corners, boxrule=0.5pt,
        colframe=#3,
        colback=#3!10,
        coltext=black,                     % 只用 tcolorbox 的原生属性控制文字颜色
        left=1em, right=1em, top=0.5em, bottom=0.5em,
        before skip=10pt, after skip=10pt,
        separator sign none
    }
    
    % 4. 同样修复不带星号的无编号版本
    \tcolorboxenvironment{#1*}{
        standard, breakable, rounded corners, boxrule=0.5pt,
        colframe=#3,
        colback=#3!10,
        coltext=black,                     % 必须加上
        left=1em, right=1em, top=0.5em, bottom=0.5em,
        before skip=10pt, after skip=10pt,
        separator sign none
    }
}


\createobsbox{theorem}{定理}{ObsTheorem}{numberwithin=section}
\createobsbox{proposition}{命题}{ObsProp}{sibling=theorem}
\createobsbox{lemma}{引理}{ObsLemma}{sibling=theorem}
\createobsbox{corollary}{推论}{ObsCor}{sibling=theorem}
\createobsbox{claim}{断言}{ObsClaim}{sibling=theorem}
\createobsbox{definition}{定义}{ObsDef}{numberwithin=section}
\createobsbox{problem}{问题}{ObsPro}{numberwithin=section}
\createobsbox{exercise}{习题}{ObsEx}{numberwithin=section}
\createobsbox{example}{例题}{ObsEg}{numberwithin=section}
\createobsbox{remark}{注记}{ObsRemark}{numberwithin=section}
\createobsbox{note}{注意}{ObsNote}{numberwithin=section}
% Solution 环境单独定义 (解答一般不换行，且需要结束符)
\declaretheoremstyle[
    headfont=\bfseries\color{ObsSol},
    bodyfont=\normalfont,
    qed=\ensuremath{\square},
    headpunct={.},
    postheadspace=1em
]{solstyle}
\declaretheorem[style=solstyle, name=Solution, numbered=no]{solution}
\tcolorboxenvironment{solution}{
    standard, breakable, rounded corners, boxrule=0.5pt,
    colframe=ObsSol, colback=ObsSol!10,
    left=1em, right=1em, top=0.5em, bottom=0.5em,
    before skip=10pt, after skip=10pt
}

\titleformat{\section}{\centering\normalfont\Large\bfseries}{\thesection}{1em}{}